\section{Architecture}\label{sec:architecture}

\begin{figure}[t]
  \centering
  \includegraphics[width=\linewidth]{mossvl-framework.pdf}
  \caption{\mossvl{} architecture. Images and video frames are encoded by
  a 27-layer vision encoder at native dynamic resolution, then pooled
  $2\times2$ and projected into visual tokens; text is tokenized alongside.
  The LLM decoder (48 layers, initialized from Qwen3-8B) attends to visual
  tokens only through 12 tanh-gated cross-attention layers
  (Gated-XAttention, right inset), one in every four; the other 36
  self-attention layers operate on the text sequence alone, so visual
  tokens never join the decoded sequence. Cross-attention queries carry
  text positions and keys carry three-axis XRoPE coordinates $(t,h,w)$.
  Gates are zero-initialized scalars, so training starts from an intact
  language backbone.}
  \label{fig:arch}
\end{figure}

\begin{figure}[t]
  \centering
  \includegraphics[width=\linewidth]{xrope.pdf}
  \caption{XRoPE, the position encoding of the cross-attention channel.
  Text tokens and visual patches share one three-axis coordinate space
  $(t,h,w)$, ordered by their logical position in the stream: a text
  token advances all three axes together, while a frame anchored at
  coordinate $t$ tiles its patches from $(t,t,t)$ at the top-left to
  $(t,t{+}h{-}1,t{+}w{-}1)$ at the bottom-right. The separator token
  closing each frame on the vision side and that frame's placeholder
  token in the text stream receive the same coordinate, so both channels
  advance along a single shared timeline. Rotations are applied to
  text-side queries and vision-side keys before they meet in
  cross-attention.}
  \label{fig:xrope}
\end{figure}

\begin{table}[t]
  \centering
  \caption{\mossvl{} configuration, identical across all released
  checkpoints. Video-input ranges reflect the native training setup; the
  released processor defaults to 1\,fps and 256 frames, and our
  evaluations keep 1\,fps with the cap raised to 768 frames
  (\S\ref{sec:evaluation}).}
  \label{tab:config}
  \small
\setlength{\tabcolsep}{6pt}
\renewcommand{\arraystretch}{1.15}
\begin{tabular}{@{}ll@{}}
\toprule
\multicolumn{2}{@{}l}{\itshape Language decoder (initialized from Qwen3-8B)} \\
\midrule
Layers            & 48 = 36 self-attention + 12 gated cross-attention \\
Cross-attention placement & every 4th layer (indices $2, 6, \ldots, 46$) \\
Hidden size / FFN size    & 4096 / 12288 \\
Attention heads   & 32 query / 8 key--value (GQA), head dim.\ 128 \\
Position encoding & XRoPE: interleaved 3-axis RoPE $(t,h,w)$, \\
                  & \quad sections $(24,20,20)$, base $5\times10^{6}$ \\
Context window    & 262{,}144 tokens \\
Vocabulary        & 151{,}936 \\
\midrule
\multicolumn{2}{@{}l}{\itshape Vision encoder (initialized from Qwen3-VL)} \\
\midrule
Layers            & 27 \\
Hidden size / FFN size & 1152 / 4304 \\
Attention heads   & 16 \\
Patch size        & $16\times16$ spatial, 1 frame temporal \\
Feature levels    & layers $\{8, 16, 24\}$ + final layer \\
Projection        & $2\times2$ spatial merge + MLP $\rightarrow$ 4096 \\
Input resolution  & native dynamic, 4096--16.8M pixels \\
\midrule
\multicolumn{2}{@{}l}{\itshape Video input} \\
\midrule
Frame sampling    & dynamic 1--16 fps, motion-adaptive (1--2 fps typical) \\
Frames per video  & up to 2{,}048 in training; 768 in our evaluations
                    (released default 256) \\
\midrule
Total parameters  & 11.3B (BF16) \\
\bottomrule
\end{tabular}

\end{table}

\begin{table}[t]
  \centering
  \caption{Released \mossvl{} checkpoints. The 0708 run is the subject of
  this report: \mossbase{} is the pre-trained model, \mossinstruct{} is its
  instruction-tuned successor, and \mossrealtime{} continues training from
  \mossinstruct{} with Realtime-SFT. The 0408 checkpoints are an earlier,
  independently trained run of the same architecture, released alongside
  for research continuity. All five are available under the OpenMOSS-Team
  organization on HuggingFace.}
  \label{tab:lineage}
  \small
\setlength{\tabcolsep}{5pt}
\renewcommand{\arraystretch}{1.15}
\begin{tabular}{@{}llll@{}}
\toprule
Checkpoint & Initialized from & Training & Evaluated in \\
\midrule
MOSS-VL-Base-0708     & --- & pre-training (\S\ref{sec:pretraining}) & --- \\
MOSS-VL-Instruct-0708 & Base-0708 & SFT (\S\ref{sec:posttraining}) & offline (\Cref{tab:offline}) \\
MOSS-VL-Realtime      & Instruct-0708 & Realtime-SFT (\S\ref{sec:posttraining}) & streaming (\Cref{tab:streaming}) \\
\addlinespace
MOSS-VL-Base-0408     & --- & pre-training & --- \\
MOSS-VL-Instruct-0408 & Base-0408 & SFT & --- \\
\bottomrule
\end{tabular}

\end{table}


\mossvl{} pairs a native-resolution vision encoder with a language
decoder initialized from Qwen3-8B \citep{Yang2025qwen}, and the two
interact only through gated cross-attention (Figure~\ref{fig:arch}).
Visual tokens never enter the decoded sequence: each frame contributes
a few timestamp tokens and one placeholder token to the text stream,
while its patch tokens are consumed as cross-attention keys and values.
Table~\ref{tab:config} lists the configuration.

\subsection{Components and Parameter Accounting}\label{sec:arch-components}

The vision encoder is a 27-layer transformer initialized from the
Qwen3-VL vision encoder \citep{Bai2025qwen}; it processes images and
frames at native resolution, from 4{,}096 to 16.8M pixels, drawing
features from three intermediate layers and the final layer. The
projection module merges each $2\times2$ patch group and maps it into
the decoder's hidden space. The decoder stacks 48 layers: 36
self-attention layers carried over from Qwen3-8B, and 12 gated
cross-attention layers, one at every fourth position. Of the 11.3B
total parameters, roughly 8.2B form the language backbone, 2.3B the
cross-attention stack, and 0.8B the vision encoder and projection
module.

\subsection{Gated Cross-Attention}\label{sec:xattn}

Each cross-attention layer follows the gated design of Flamingo
\citep{Alayrac2022flamingo}---queries come from the text hidden
states, keys and values from the visual tokens---implemented here
with grouped-query attention (32 query / 8 key--value heads) and
QK-RMSNorm. The layer wraps its
attention and feed-forward paths in tanh gates whose scalars are
zero-initialized, so training starts from an intact language backbone
(\S\ref{sec:pretraining}). The 36 self-attention layers never see
visual tokens.

\subsection{XRoPE}\label{sec:xrope}

To our knowledge, XRoPE (cross-attention rotary position embedding) is
the first position encoding introduced for the cross-attention channel
of a vision--language architecture: it gives the visual stream position
information that this channel otherwise lacks. XRoPE places text tokens
and visual patches in one three-axis coordinate space $(t,h,w)$,
ordered by their logical position in the stream
(Figure~\ref{fig:xrope}). A text token advances all three axes
together, taking coordinate $(x,x,x)$. A frame whose merged patch grid
is $h'\times w'$ anchors at the coordinate $t$ following the preceding
text, and its patches tile
\begin{equation}\label{eq:xrope}
  \mathbf{p}_{a,b} = (t,\; t+a,\; t+b),
  \qquad 0 \le a < h',\;\; 0 \le b < w',
\end{equation}
so height and width offsets ride on the shared temporal anchor. The
separator token that closes the frame on the vision side and the
frame's placeholder token in the text stream both take $(x,x,x)$ with
$x = \max(t+h',\, t+w')$, and the next text token continues from
$x+1$: the two channels advance along a single timeline. The 64 rotary
frequency pairs are split $(24,20,20)$ across $(t,h,w)$, and rotations
are applied to text-side queries and vision-side keys before they meet
in cross-attention. The $t$ axis is a relative sequence coordinate,
not wall-clock time; real timing enters through the timestamp tokens
of \S\ref{sec:timestamps}.

\subsection{Absolute Timestamps}\label{sec:timestamps}

Positions alone say nothing about wall-clock time, and frame rates
vary: \mossvl{} samples video at 1--16\,fps, motion-adaptive
(Table~\ref{tab:config}). Each frame is therefore preceded in the text
stream by an absolute timestamp,
\statetok{time\_start}\texttt{X.X seconds}\statetok{time\_end}, so
the model reads real time from tokens rather than inferring it from
positions, and timing stays explicit under any sampling rate.

\subsection{Real-Time by Construction}\label{sec:realtime-arch}

When a new frame arrives, only that frame is encoded; its keys and
values are appended to the cross-attention cache, and earlier frames
are neither re-encoded nor their keys and values recomputed. The
decoded sequence grows by the frame's
timestamp tokens and a single placeholder token---patch tokens stay on
the vision side---so an arriving stream leaves the decoding state
intact, and the next generated token already attends to the updated
cache through the gated layers.
\S\ref{sec:evaluation} quantifies the efficiency this
yields at inference time (Figure~\ref{fig:efficiency}).

\subsection{Released Models}\label{sec:released-models}

We release five checkpoints of one architecture
(Table~\ref{tab:lineage}). The 0708 run---\mossbase{},
\mossinstruct{}, \mossrealtime{}---is the subject of this report; the
0408 pair is an earlier, independently trained run of the same
architecture, released for research continuity.
